\makeabstract
    {
    Brane trajectories in cosmological spacetimes
    }
    {
    Yunwen Sophia Xu\textsuperscript{1}, David Kastor\textsuperscript{2}
    }
    {
    \textsuperscript{1} Department of Physics \& Astronomy, Amherst College, Amherst, MA\\
\textsuperscript{2} Department of Physics, University of Massachusetts at Amherst, Amherst, MA
    }
    {
    A fundamental mystery of our universe is its particular rate of expansion, fueled by an outward pressure known as dark energy. While we can observe the impacts of dark energy on our universe, we cannot experimentally compare different universes to determine the causes or consequences of a specific dark energy content. Instead, we use theoretical methods to better understand the role of dark energy in cosmology through describing the behavior of objects in 3+1-dimensional spacetimes.
    }
    {
    We apply mathematical tools from general relativity, differential geometry, and string theory to study the trajectories of general-dimensional objects moving in spacetimes, known as branes. By applying the variational principle to brane area-action, we derive the equation of motion for a general brane which is dependent on its geometry. We then apply this formalism to compare the trajectories of massive, spherical 2-branes in Minkowski and de Sitter spacetimes. 
    }
    {
    The equations of motion for a massive, spherical 2-brane in Minkowski spacetime demonstrate that it tends toward shrinking, regardless of initial conditions. By contrast, a massive, spherical 2-brane in de Sitter spacetime has a critical radius at which the inward surface tension and outward pressure perfectly balance. When the brane is larger than the critical radius, it tends toward expanding as the outward force due to dark energy is greater than inward force from surface tension. Correspondingly, when the brane is smaller than the critical radius, it tends toward shrinking. These results align with their classical analogue, where higher pressure within a bubble allows it to hold its form, while no pressure difference would lead to a shrinking bubble.
    }
    {
    While the dark energy content of a spacetime is conventionally seen as a constant, branes with negative energy can be absorbed by black holes, leading to a change in the dark energy content of the spacetime. This research opens the door to further study of mechanisms for changing the dark energy content of spacetimes via black hole energy extraction. Our eventual goal is to describe negative-energy charged branes absorbed by a rotating Ba\~nados-Teitelboim-Zanelli (BTZ) black hole, as a model mechanism enabling the fine-tuning of our universe's dark energy content.
    }

    \newpage
    